A2Perf is a benchmark harness designed to be used flexibly on various machines. The user has the option to either run it in a Docker container or to run the benchmark locally. A Docker container is available in the harness and can be adapted to your needs. If you would like to run the benchmark locally, a guide is available to install the A2Perf benchmark harness on your Linux or MacOS system. While you can benchmark on both operating systems, it is important to note that system performance metrics are tracked using CodeCarbon. This allows to capture energy, power and memory usage at regular time intervals, and uses pyRAPL to compute the Running Average Power Limit (RAPL). However, RAPL uniquely measures power consumption information for Intel CPUs, DRAM for server architectures and GPU for client architectures. When using systems using CPU architectures different then the Intel CPUs, the power consumption metric will return a computed estimate rather than a measured metric.

The benchmark harness allows you to benchmark both the training and inference of your algorithm and agents respectively. In order to benchmark your algorithms, you need to create a submission folder which includes several files which A2Perf calls. First, a training file, \texttt{train.py} contains a function \texttt{train()}, which starts the training process of your algorithm when called. Similarly, \texttt{inference.py} covers the inference of your trained model. This file includes several functions responsible for the loading of your trained model, preprocessing observations and running inference on your model. Using a \texttt{requirements.txt} file, additional Python packages and versioning can be specified. Running the benchmark is done through a command line interface. Using flags, we can pass additional information to the submission to set up the benchmark. A \texttt{gin\_config} flag allows the user to define the settings for your environment and training process. Additionally, we need to pass the path to the submission folder using the \texttt{participant\_module\_path} flag. For a more detailed description, tutorials are available in the repository.
